\chapter{Marco metodológico}\label{chap:marco}

Cuando se inició el proyecto se descartó la idea de diseñar una solución completa antes de escribir código porque el problema de partida, extraer datos de carátulas bancarias con formatos heterogéneos, presentaba demasiadas incógnitas como para definir una arquitectura definitiva sin antes probarla con documentos reales. Por ello se adoptó un enfoque iterativo inspirado en \textit{Lean Startup} \cite{ries_lean_startup}: construir algo mínimo, medir su desempeño y usar los resultados para decidir el siguiente paso. Se inició experimentando con enfoques determinísticos: expresiones regulares sobre texto extraído con \textit{pdfplumber}, combinaciones con \textit{OCR} y variantes similares. Los resultados fueron contundentes: ninguna de estas estrategias superó el 13\% de precisión promedio, lo que dejó claro desde temprano que los métodos basados en patrones fijos no iban a funcionar con documentos con alta variabilidad estructural. Por ello, se exploraron variantes con modelos de lenguaje (Claude con texto extraído, con imágenes, y con \textit{OCR} como respaldo), y la combinación de \textit{OCR} con Claude mostró los mejores resultados de esa ronda, alcanzando un 63.8\% de precisión promedio. Con esos datos, se construyó el prototipo funcional, descrito con mayor detalle en el capítulo~\ref{chap:prototipo}, usando precisamente ese extractor híbrido de \textit{OCR} + Claude, con el objetivo de responder una pregunta concreta: ¿puede este enfoque sostener un flujo \textit{end-to-end} sobre documentos reales? Se probó en cinco documentos y el prototipo funcionó, pero expuso limitaciones. El \textit{OCR} introducía errores, el campo Banco aparecía frecuentemente como logotipo y no como texto reconocible, mientras que la variabilidad entre formatos bancarios era mayor de lo esperado. Además, la API de Claude incorporó nuevas capacidades para procesar PDFs directamente como binario, lo que motivó una tercera ronda de experimentación (capítulo~\ref{chap:experimentos}, fase~3). Esta vez se evaluó el envío directo del PDF al modelo sin ningún preprocesamiento intermedio, y los resultados fueron decisivos: 93.4\% de precisión promedio, superando ampliamente al extractor del prototipo. La mejora más notable fue en el campo Banco, que pasó de 41.5\% a 96.6\% al eliminar la dependencia de \textit{OCR} para interpretar logotipos.

Finalmente, con la arquitectura de extracción validada, se construyó el producto mínimo viable (capítulo~\ref{chap:mvp}). Aquí el desafío ya no era técnico sino de producto: ¿cómo hacer que esto sea útil para una \textit{fintech} real? La variabilidad de precisión entre bancos (0\% en BanBajío, 100\% en HSBC) descartó la idea de un extractor único y universal, lo que llevó a diseñar un sistema de extractores configurables donde cada usuario puede definir su propio \textit{schema}, \textit{prompt} y modelo. Esta decisión de arquitectura no habría surgido sin los datos de las etapas experimentales. La Figura~\ref{fig:marco_metodologico} resume esta secuencia y las hipótesis que cada etapa buscó validar. En retrospectiva, lo que hizo funcionar este enfoque fue no comprometerse con una arquitectura antes de tener datos. Cada etapa se justificó con los resultados de la anterior, y varias decisiones importantes, como abandonar el \textit{OCR} o diseñar extractores configurables, surgieron de hallazgos que no eran predecibles al inicio del proyecto.
\begin{figure}[H]
\centering
\resizebox{\textwidth}{!}{
\begin{tikzpicture}[
    node distance=0.6cm and 0.9cm,
    every node/.style={font=\small, align=center},
    etapa/.style={
        rectangle,
        rounded corners=6pt,
        draw=black,
        thick,
        minimum width=3.2cm,
        minimum height=2.2cm,
        fill=cyan!10
    },
    aprendizaje/.style={
        font=\scriptsize\itshape,
        text width=2.6cm,
        align=center
    },
    flecha/.style={
        -{Latex[length=3mm]},
        thick
    },
    ciclo/.style={
        -{Latex[length=2.5mm]},
        thick,
        gray!60,
        dashed
    }
]

% Etapas
\node[etapa] (exp1) {
    Experimentación\\inicial\\[3pt]
    {\scriptsize\textit{Cap.~\ref{chap:experimentos}, fases 1--2}}
};
\node[etapa, right=2.2cm of exp1] (proto) {
    Prototipo\\funcional\\[3pt]
    {\scriptsize\textit{Cap.~\ref{chap:prototipo}}}
};
\node[etapa, right=2.2cm of proto] (exp2) {
    Experimentación\\PDF directo\\[3pt]
    {\scriptsize\textit{Cap.~\ref{chap:experimentos}, fase 3}}
};
\node[etapa, right=2.2cm of exp2] (mvp) {
    Producto mínimo\\viable (MVP)\\[3pt]
    {\scriptsize\textit{Cap.~\ref{chap:mvp}}}
};

% Flechas entre etapas con aprendizajes
\draw[flecha] (exp1) -- node[above, aprendizaje] {\textit{Regex} falla;\\OCR + Claude\\es viable} (proto);
\draw[flecha] (proto) -- node[above, aprendizaje] {\textit{OCR} introduce\\errores; Banco\\sale como logo} (exp2);
\draw[flecha] (exp2) -- node[above, aprendizaje] {PDF directo\\gana; precisión\\varía por banco} (mvp);

% Hipótesis debajo
\node[below=0.4cm of exp1, text width=3.2cm, align=center, font=\scriptsize] {
    Hipótesis: \textit{LLMs} superan\\a métodos determinísticos
};
\node[below=0.4cm of proto, text width=3.2cm, align=center, font=\scriptsize] {
    Hipótesis: El flujo\\funciona \textit{end-to-end}
};
\node[below=0.4cm of exp2, text width=3.2cm, align=center, font=\scriptsize] {
    Hipótesis: Simplificar el\\\textit{pipeline} mejora precisión
};
\node[below=0.4cm of mvp, text width=3.2cm, align=center, font=\scriptsize] {
    Hipótesis: Extractores\\configurables escalan
};

% Ciclo construir-medir-aprender
\draw[ciclo] (mvp.south) .. controls +(0,-2.2) and +(-0.5,-2.2) .. node[below, font=\scriptsize, text=gray!70] {Construir — Medir — Aprender} (exp1.south);

\end{tikzpicture}
}
\caption{Enfoque metodológico iterativo: experimentación, prototipado, experimentación y producto.}
\label{fig:marco_metodologico}
\end{figure}


